\documentclass{beamer}
	\usetheme{Dresden}
	\usepackage{beamerthemesplit}
	
	\useoutertheme{split}
	\usecolortheme{whale}
	\useinnertheme{rounded} 
	\usecolortheme{orchid}
	\setbeamertemplate{blocks}[rounded]


\usepackage{listings}  % For code formatting
\usepackage{xcolor}    % For color support in code

% Configure listings
\lstset{
	basicstyle=\ttfamily\tiny, 
	keywordstyle=\color{blue}\bfseries,
	commentstyle=\color{green},
	stringstyle=\color{red},
 	numberstyle=\tiny\color{codegray}, % hb
	frame=single,
	breaklines=true,         % Line breaking
	showstringspaces=false,  % Don't show spaces in strings
	tabsize=2,               % Set tab size to 4 spaces
	language=Java,           % Set language to Java for syntax highlighting
	literate={->}{$\rightarrow$}{2} % Correctly display -> as an arrow
}

% Add page numbers with total frame count in the footer
\setbeamertemplate{footline}[frame number]{}
\setbeamertemplate{footline}{%
    \hfill \insertframenumber/\inserttotalframenumber\hspace{1em}
}

% Reduce space between the slide title and the content globally
\addtobeamertemplate{frametitle}{}{\vspace{-1em}}

\title{Introduction to JUnit in Java}
\subtitle{A Guide to Unit Testing with JUnit}
\author{[Haluk Bingol]}
\date{\today}

\begin{document}

\frame{\titlepage}

% Table of Contents
\begin{frame}{Table of Contents}
    \tableofcontents
\end{frame}


\section{What is JUnit?}

\begin{frame}{What is JUnit?}
	\begin{itemize}
		\item JUnit is a popular testing framework for Java.
		\item Supports writing and running tests to ensure code quality.
		\item Follows the TDD (Test-Driven Development) approach.
		\item Integrated with most Java IDEs (e.g., Eclipse, IntelliJ IDEA).
	\end{itemize}
\end{frame}

\begin{frame}{Why Use JUnit?}
	\begin{itemize}
		\item Automated Testing: Efficient and repeatable tests.
		\item Early Bug Detection: Helps catch bugs early in development.
		\item Regression Testing: Ensures new code doesn’t break existing functionality.
		\item Improves Code Quality: Encourages modular, testable code.
	\end{itemize}
\end{frame}

\section{JUnit Architecture}

\begin{frame}{JUnit Architecture}
	\begin{itemize}
		\item \textbf{Test Runner}: Executes tests.
		\item \textbf{Test Fixtures}: Prepare setup and teardown routines.
		\item \textbf{Test Cases}: Individual units of tests (methods).
		\item \textbf{Assertions}: Validate expected outcomes.
		\item \textbf{Annotations}: Simplify test configurations (@Test, @Before, etc.).
	\end{itemize}
\end{frame}

\begin{frame}{Key Annotations in JUnit}
	\begin{itemize}
		\item \textbf{@Test}: Marks a method as a test method.
		\item \textbf{@BeforeEach}: Executes before each test (setup).
		\item \textbf{@AfterEach}: Executes after each test (cleanup).
		\item \textbf{@BeforeAll}: Runs once before all tests.
		\item \textbf{@AfterAll}: Runs once after all tests.
		\item \textbf{@Disabled}: Disables a test temporarily.
	\end{itemize}
\end{frame}

\section{a Simple Test Case}

\begin{frame}[fragile]{Creating a Simple Test Case}
	Example:
	\begin{lstlisting}
import static org.junit.jupiter.api.Assertions.*;
import org.junit.jupiter.api.Test;

public class CalculatorTest {
    @Test
    void testAdd() {
        Calculator calc = new Calculator();
        assertEquals(5, calc.add(2, 3), "2 + 3 should be 5");
    }
}
	\end{lstlisting}
	\textbf{assertEquals}: Checks if the actual value matches the expected value.
\end{frame}

\begin{frame}{Common Assertions in JUnit}
	\begin{itemize}
		\item \textbf{assertEquals(expected, actual)}: Asserts equality.
		\item \textbf{assertTrue(condition)}: Asserts if the condition is true.
		\item \textbf{assertFalse(condition)}: Asserts if the condition is false.
		\item \textbf{assertNull(object)}: Asserts if the object is null.
		\item \textbf{assertNotNull(object)}: Asserts if the object is not null.
		\item \textbf{assertThrows(Exception.class, () $\rightarrow$ methodCall)}: Asserts if an exception is thrown.
	\end{itemize}
\end{frame}

\subsection{Parameterized Tests in JUnit}

\begin{frame}[fragile]{Parameterized Tests in JUnit}
	Parameterized tests allow testing with multiple sets of data.

	Example:
	\begin{lstlisting}
@ParameterizedTest
@ValueSource(ints = {2, 4, 6})
void testIsEven(int number) {
    assertTrue(number % 2 == 0);
}
	\end{lstlisting}
	This test will run three times, each time with different input.
\end{frame}

\subsection{JUnit Test Lifecycle}

\begin{frame}{JUnit Test Lifecycle}
	\begin{itemize}
		\item \textbf{@BeforeAll}: Set up shared resources.
		\item \textbf{@BeforeEach}: Set up before each test (e.g., initialize objects).
		\item \textbf{@Test}: Execute the test case.
		\item \textbf{@AfterEach}: Tear down (cleanup) after each test.
		\item \textbf{@AfterAll}: Clean up shared resources.
	\end{itemize}
\end{frame}

\subsection{Handling Exceptions }

\begin{frame}[fragile]{Handling Exceptions in JUnit}
	Example of testing exception handling:
	\begin{lstlisting}
@Test
void testDivisionByZero() {
    assertThrows(ArithmeticException.class, () -> {
        Calculator calc = new Calculator();
        calc.divide(1, 0);
    });
}
	\end{lstlisting}
	Verifies that an exception is correctly thrown under specific conditions.
\end{frame}

\begin{frame}[fragile]{Mocking with JUnit}
	\begin{itemize}
		\item Mockito can be used with JUnit to mock dependencies.
		\item Example:
	\end{itemize}
	\begin{lstlisting}
@Mock
List<String> mockedList;

@Test
void testMockList() {
    when(mockedList.size()).thenReturn(5);
    assertEquals(5, mockedList.size());
}
	\end{lstlisting}
\end{frame}

\subsection{Best Practices}

\begin{frame}{Best Practices for JUnit Testing}
	\begin{itemize}
		\item Write \textbf{independent} test cases.
		\item Name tests clearly to reflect their purpose.
		\item Use \textbf{mocking} for external dependencies.
		\item Keep tests \textbf{simple and focused} on one behavior.
		\item Aim for \textbf{100\% code coverage} where feasible.
	\end{itemize}
\end{frame}

\begin{frame}{JUnit 5 vs. JUnit 4}
	\begin{itemize}
		\item JUnit 5 introduced:
		\begin{itemize}
			\item Modular architecture (JUnit Platform, JUnit Jupiter, JUnit Vintage).
			\item \textbf{@BeforeEach} and \textbf{@AfterEach} replace \textbf{@Before} and \textbf{@After}.
			\item More powerful extensions and parameterized tests.
			\item Support for Java 8 and higher (lambda expressions).
		\end{itemize}
	\end{itemize}
\end{frame}

\subsection{Running JUnit Tests in IDEs}

\begin{frame}{Running JUnit Tests in IDEs}
	\begin{itemize}
		\item \textbf{Eclipse/IntelliJ}: Right-click the test class or method and select "Run as JUnit Test."
		\item \textbf{Maven/Gradle Integration}: JUnit tests can be run as part of the build process.
		\item Example Maven command: \texttt{mvn test}
	\end{itemize}
\end{frame}

\begin{frame}{Conclusion}
	\begin{itemize}
		\item JUnit is a vital tool for Java developers.
		\item Helps maintain code quality and catch bugs early.
		\item Supports automated testing with flexible configurations.
		\item Integrates with modern development tools and methodologies (TDD, CI/CD).
	\end{itemize}
\end{frame}

\begin{frame}{Questions?}
	\centering
	Any questions?
\end{frame}

\end{document}
